\documentclass[12pt]{article}
\usepackage{amsmath}
\usepackage{amsfonts}
\usepackage{amssymb}
\usepackage{physics}
\usepackage{geometry}
\usepackage{parskip}
\geometry{a4paper, margin=0.7in}


\title{02YMECH - Solution 13}
\date{Assigned for the week of Dec 15, 2025}
\author{}

\begin{document}
\maketitle

\section*{Solution}

(a) For torque-free motion, Euler's equations are
\[
\begin{aligned}
I_1\dot{\omega}_1 &= (I_2-I_3)\omega_2\omega_3,\\
I_2\dot{\omega}_2 &= (I_3-I_1)\omega_3\omega_1,\\
I_3\dot{\omega}_3 &= (I_1-I_2)\omega_1\omega_2.
\end{aligned}
\]

The kinetic energy is
\[
T=\frac12\left(I_1\omega_1^2+I_2\omega_2^2+I_3\omega_3^2\right).
\]
Differentiating and using Euler’s equations gives
\[
\dot T = I_1\omega_1\dot{\omega}_1 + I_2\omega_2\dot{\omega}_2 + I_3\omega_3\dot{\omega}_3,
\]
and substituting Euler's equations yields
\[
\dot T = \omega_1(I_2-I_3)\omega_2\omega_3 + \omega_2(I_3-I_1)\omega_3\omega_1 + \omega_3(I_1-I_2)\omega_1\omega_2 = 0,
\]
\[
\dot T=0,
\]
so \(T\) is conserved.

The angular momentum in the body frame is
\[
\mathbf L=(I_1\omega_1,I_2\omega_2,I_3\omega_3),
\]
and its squared magnitude is
\[
L^2=I_1^2\omega_1^2+I_2^2\omega_2^2+I_3^2\omega_3^2.
\]
Differentiating,
\[
\dot L^2 = 2I_1^2\omega_1\dot{\omega}_1 + 2I_2^2\omega_2\dot{\omega}_2 + 2I_3^2\omega_3\dot{\omega}_3,
\]
and substituting Euler's equations yields
\[
\dot L^2 = 2I_1\omega_1(I_2-I_3)\omega_2\omega_3 + 2I_2\omega_2(I_3-I_1)\omega_3\omega_1 + 2I_3\omega_3(I_1-I_2)\omega_1\omega_2 = 0,
\]
satisfies \(\dot L^2=0\). Hence \(L^2\) is conserved. \emph{Even though \(\mathbf L\) itself is not conserved in the body frame, its magnitude is.}


(b) From the initial conditions,
\[
T=\frac12(I_1+I_2)\omega_0^2,
\qquad
L^2=(I_1^2+I_2^2)\omega_0^2.
\]

Thus, for all \(t\),
\begin{align}
I_1\omega_1^2+I_2\omega_2^2+I_3\omega_3^2
&=(I_1+I_2)\omega_0^2, \label{E}\\
I_1^2\omega_1^2+I_2^2\omega_2^2+I_3^2\omega_3^2
&=(I_1^2+I_2^2)\omega_0^2. \label{L}
\end{align}

Solving \eqref{E}--\eqref{L} for \(\omega_2^2\) and \(\omega_3^2\) yields
\[
\omega_3^2
=
\frac{I_2-I_1}{I_3-I_1}\,(\omega_1^2-\omega_0^2),
\qquad
\omega_2^2
=
\frac{I_1(I_3-I_2)}{I_2(I_2-I_1)}(\omega_0^2-\omega_1^2).
\]

From the first Euler equation,
\[
\dot{\omega}_1=\frac{I_2-I_3}{I_1}\omega_2\omega_3.
\]
Squaring and substituting the expressions above gives
\[
\boxed{
\dot{\omega}_1^{\,2}
=
\frac{(I_3-I_2)(I_2-I_1)}{I_1^2(I_3-I_1)}
\,\omega_1^2(\omega_0^2-\omega_1^2)
}.
\]



(c) Since \(\dot{\omega}_1^{\,2}\ge0\),
\[
0\le \omega_1^2\le \omega_0^2,
\quad\text{or}\quad
-\omega_0\le\omega_1\le\omega_0.
\]
The motion is bounded between turning points and therefore periodic.


(d) Let
\[
C=\frac{(I_3-I_2)(I_2-I_1)}{I_1^2(I_3-I_1)}.
\]
Then
\[
\frac{d\omega_1}{dt}
=
\sqrt{C}\,\omega_1\sqrt{\omega_0^2-\omega_1^2}.
\]
Separating variables,
\[
dt
=
\frac{d\omega_1}{\sqrt{C}\,\omega_1\sqrt{\omega_0^2-\omega_1^2}}.
\]

One quarter of the period corresponds to \(\omega_1:0\to\omega_0\):
\[
\frac{T}{4}
=
\frac{1}{\sqrt{C}}
\int_0^{\omega_0}
\frac{d\omega_1}{\omega_1\sqrt{\omega_0^2-\omega_1^2}}.
\]

With the substitution \(\omega_1=\omega_0\sin\theta\),
\[
\boxed{
T
=
\frac{4}{\omega_0\sqrt{C}}
\int_0^{\pi/2}\frac{d\theta}{\sin\theta}
},
\]
which is an integral representation of the motion.
\end{document}
